% =========================================================================
%  Supplementary Materials for:
%  Finite Prevention Windows for HIV Post-Exposure Prophylaxis:
%  Irreversible Proviral Integration Defines Route-Specific Intervention Limits
%
%  A.C. Demidont, Nyx Dynamics LLC
%  v2 supplement, companion to manuscript v2 for Science Advances aeh5879
% =========================================================================
\documentclass[11pt]{article}

\usepackage[T1]{fontenc}
\usepackage[utf8]{inputenc}
\usepackage{mathptmx}
\usepackage[scaled]{helvet}

\usepackage[letterpaper,margin=1in]{geometry}
\usepackage{setspace}
\onehalfspacing

\usepackage{amsmath,amssymb,amsthm}
\usepackage{bm}
\usepackage{mathtools}

\usepackage{graphicx}
\usepackage{booktabs}
\usepackage{array}
\usepackage{tabularx}
\usepackage{longtable}
\usepackage[font=small,labelfont=bf]{caption}
\usepackage{float}

\usepackage{enumitem}
\usepackage{xcolor}
\usepackage[colorlinks=true,linkcolor=black,citecolor=black,urlcolor=blue]{hyperref}

% Theorem environments — SI numbering
\newtheorem{theorem}{Theorem}[section]
\newtheorem{lemma}{Lemma}[section]
\newtheorem{proposition}{Proposition}[section]
\newtheorem{corollary}{Corollary}[section]
\theoremstyle{definition}
\newtheorem{definition}{Definition}[section]
\newtheorem{assumption}{Assumption}[section]
\theoremstyle{remark}
\newtheorem*{remark}{Remark}

% Macros (mirror main file)
\newcommand{\Tseed}{T_{\text{seed}}}
\newcommand{\Tint}{T_{\text{int}}}
\newcommand{\Pseed}{P_{\text{seed}}}
\newcommand{\Pint}{P_{\text{int}}}
\newcommand{\Epep}{E_{\text{PEP}}}
\newcommand{\Ebarpep}{\bar{E}_{\text{PEP}}}
\newcommand{\Faccess}{F_{\text{access}}}
\newcommand{\tcrit}{t_{\text{crit}}}
\newcommand{\emax}{\varepsilon_{\max}}
\newcommand{\emid}{\varepsilon_{\text{mid}}}
\newcommand{\emin}{\varepsilon_{\min}}

\setlength{\parskip}{0.4em}
\setlength{\parindent}{0pt}

% S-prefixed numbering for Science-style supplement
\renewcommand{\thesection}{S\arabic{section}}
\renewcommand{\thesubsection}{S\arabic{section}.\arabic{subsection}}
\renewcommand{\thefigure}{S\arabic{figure}}
\renewcommand{\thetable}{S\arabic{table}}
\renewcommand{\theequation}{S\arabic{equation}}

\usepackage{titlesec}
\titleformat*{\section}{\large\bfseries}
\titleformat*{\subsection}{\normalsize\bfseries}

% =========================================================================
\begin{document}

\begin{flushleft}
{\Large\bfseries Supplementary Materials for\par}
\vspace{0.3em}
{\large\bfseries Finite Prevention Windows for HIV Post-Exposure Prophylaxis: Irreversible Proviral Integration Defines Route-Specific Intervention Limits\par}
\vspace{0.8em}
A.C. Demidont\textsuperscript{*}\\
\textit{Nyx Dynamics LLC, Philadelphia, PA 19107, USA}\\
\textsuperscript{*}Correspondence: \texttt{ac.demidont@nyxdynamics.com}\\
ORCID: 0000-0002-9216-8569
\end{flushleft}

\vspace{0.5em}
\hrule
\vspace{0.5em}

\noindent\textbf{This PDF file includes:}
\begin{itemize}[leftmargin=2em,nosep]
    \item Sections~S1--S6: Mathematical development with complete proofs (unchanged from v1)
    \item Section~S7: Multiscale within-host model implementation
    \item Section~S8: NHP concordance --- per-timepoint detail
    \item Section~S9: Parameter tables and stochastic VL methods
    \item Section~S10: City-stratified extended methods
    \item Section~S11: Figure inventory
\end{itemize}

\vspace{0.5em}
\noindent\textbf{Note on v2 changes.} The mathematical content in Sections~S1--S6 (theorems, proofs, lemmas, propositions) is unchanged from v1. Numerical worked examples in Section~S6 have been updated to reflect verified values from the multiscale model implementation (commit \texttt{37e27ea}). New material in Sections~S7--S8 documents the multiscale implementation and NHP concordance.

% =========================================================================
\section{Probability Space, State Definitions, and Assumptions}

Consider a single HIV exposure event $e$ occurring at time $t = 0$ on probability space $(\Omega, \mathcal{F}, \mathbb{P})$ with filtration $\{\mathcal{F}_{t}\}_{t \geq 0}$.

\subsection{Within-host states}

\begin{definition}[Within-host infection states]
$Z(t) \in \{0, 1, 2\}$:
\begin{itemize}[nosep,leftmargin=2em]
    \item $Z(t) = 0$ (Susceptible): no productive infection; $I(t) < I^{*}$.
    \item $Z(t) = 1$ (Seeded): replicative focus established ($I(t) \geq I^{*}$) but reservoir below irreversibility threshold ($R(t) < R^{*}$).
    \item $Z(t) = 2$ (Integrated): $R(t) \geq R^{*}$. Absorbing state under Assumption~A1.
\end{itemize}
\end{definition}

\begin{definition}[Transition stopping times]
\[
\Tseed = \inf\{t \geq 0 : Z(t) \geq 1\}, \qquad
\Tint  = \inf\{t \geq 0 : Z(t) = 2\}.
\]
We assume $\Tseed \leq \Tint$ almost surely.
\end{definition}

\begin{definition}[Cumulative transition probabilities]
$\Pseed(t) = \mathbb{P}(\Tseed \leq t)$, $\Pint(t) = \mathbb{P}(\Tint \leq t)$. Since $\Tseed \leq \Tint$ a.s., $\Pseed(t) \geq \Pint(t)$ for all $t \geq 0$.
\end{definition}

\subsection{Assumptions}

\begin{assumption}[Irreversibility of integration]
Once $Z(t) = 2$, the state is absorbing: $\mathbb{P}(Z(s) = 0 \mid Z(t) = 2) = 0$ for all $s > t$.

\textit{Biological grounding}: Integrated provirus persists for the lifetime of the host cell and propagates to daughter cells via clonal expansion. ART suppresses replication ($I \to 0$) but $dR/dt = \alpha I \to 0$: the reservoir is static under ART but does not shrink. No approved therapy achieves eradication.
\end{assumption}

\begin{assumption}[Stage-dependent drug efficacy]
\[
\varepsilon(Z(t)) =
\begin{cases}
\emax & Z(t) = 0,\\
\emid & Z(t) = 1,\\
\emin = 0 & Z(t) = 2,
\end{cases}
\qquad
\text{with } 1 \geq \emax \geq \emid \geq 0.
\]
\end{assumption}

\begin{assumption}[Monotonicity of integration]
$\Pint(t)$ is non-decreasing, $\Pint(0) = 0$, $\lim_{t\to\infty} \Pint(t) = 1$.
\end{assumption}

\begin{assumption}[Transmission-competent exposure]
Source plasma viral load is at or above the U=U threshold ($\geq 200$~copies/mL sustained). The framework is silent on exposures from virologically suppressed sources, for whom transmission risk is functionally zero (Rodger 2019, PARTNER2; Bavinton 2018, Opposites Attract).
\end{assumption}

% =========================================================================
\section{Master Equation}

\begin{proposition}[Master equation]\label{prop:master}
Under Assumption~A2,
\begin{equation}
\Epep(t) = \big(1 - \Pseed(t)\big)\emax + \big(\Pseed(t) - \Pint(t)\big)\emid + \Pint(t)\,\emin
\label{eq:master_si}
\end{equation}
\end{proposition}

\begin{proof}
Partition the sample space at time $t$ by $Z(t)$:
\[
\mathbb{P}(Z = 0) = 1 - \Pseed,\quad
\mathbb{P}(Z = 1) = \Pseed - \Pint,\quad
\mathbb{P}(Z = 2) = \Pint.
\]
These events are mutually exclusive and exhaustive. By the law of total expectation:
\[
\Epep(t) = \mathbb{E}[\varepsilon(Z(t))] = (1 - \Pseed)\,\emax + (\Pseed - \Pint)\,\emid + \Pint\,\emin. \qedhere
\]
\end{proof}

% =========================================================================
\section{The Finite Prevention Window Theorem}

\begin{theorem}[Finite Prevention Window]\label{thm:finite}
Under Assumptions~A1--A3, $\Epep(t)$ satisfies:
\begin{enumerate}[label=\textnormal{(\roman*)},leftmargin=2.5em]
    \item $\Epep(t) \leq \big(1 - \Pint(t)\big)\,\emax$;
    \item $\Epep(t)$ is non-increasing in $t$;
    \item $\Epep(t) \to 0$ as $t \to \infty$;
    \item for any $\eta \in (0,1)$, $\tcrit(\eta) = \inf\{t : \Epep(t) < \eta\}$ is finite.
\end{enumerate}
\end{theorem}

\begin{proof}
\textit{Part (i) --- Upper bound.} Replace $\emid$ with $\emax$ in (\ref{eq:master_si}):
\[
\Epep(t) \leq (1 - \Pseed)\emax + (\Pseed - \Pint)\emax + \Pint\,\emin = (1 - \Pint)\emax + \Pint\,\emin.
\]
With $\emin = 0$: $\Epep(t) \leq (1 - \Pint(t))\emax$.

\textit{Part (ii) --- Monotone decay.} For $t_{1} < t_{2}$, let $\Delta_{S} = \Pseed(t_{2}) - \Pseed(t_{1}) \geq 0$ and $\Delta_{I} = \Pint(t_{2}) - \Pint(t_{1}) \geq 0$. Then
\[
\Epep(t_{1}) - \Epep(t_{2}) = \Delta_{S}(\emax - \emid) + \Delta_{I}(\emid - \emin) \geq 0.
\]

\textit{Part (iii) --- Convergence.} By A3, $\Pint(t) \to 1$. Since $\Tseed \leq \Tint$ a.s., $\Pseed(t) \to 1$. Limits give $\Epep(t) \to 0\cdot\emax + 0\cdot\emid + 1\cdot\emin = 0$.

\textit{Part (iv) --- Finite critical time.} By (ii), $\Epep$ is non-increasing; by (iii), $\Epep \to 0 < \eta$. Therefore $\{t : \Epep(t) < \eta\}$ is non-empty, and $\tcrit(\eta)$ is finite.
\end{proof}

% =========================================================================
\section{Hazard-Rate Decomposition of Efficacy Loss}

\begin{definition}[Hazard rates]
$\lambda_{\text{seed}}(t) = f_{\text{seed}}(t)/(1 - \Pseed(t))$, $\lambda_{\text{int}}(t) = f_{\text{int}}(t)/(\Pseed(t) - \Pint(t))$.
\end{definition}

\begin{proposition}[Decomposition]\label{prop:hazard}
Under Assumptions~A1--A3 with absolute continuity,
\begin{equation}
\frac{d\Epep}{dt} = -f_{\text{seed}}(t)\,\Delta\varepsilon_{1} - f_{\text{int}}(t)\,\Delta\varepsilon_{2} \;\leq\; 0
\label{eq:hazard_si}
\end{equation}
where $\Delta\varepsilon_{1} = \emax - \emid$ and $\Delta\varepsilon_{2} = \emid - \emin$.
\end{proposition}

\begin{proof}
Differentiating (\ref{eq:master_si}):
\[
\frac{d\Epep}{dt} = -f_{\text{seed}}\emax + (f_{\text{seed}} - f_{\text{int}})\emid + f_{\text{int}}\emin = -f_{\text{seed}}\Delta\varepsilon_{1} - f_{\text{int}}\Delta\varepsilon_{2} \leq 0. \qedhere
\]
\end{proof}

\begin{remark}
Equation~\ref{eq:hazard_si} decomposes efficacy loss into a seeding channel ($f_{\text{seed}}$, gap $\Delta\varepsilon_{1}$) and an integration channel ($f_{\text{int}}$, gap $\Delta\varepsilon_{2}$). At early times, seeding dominates; as mass accumulates in state~1, the integration channel accelerates terminal efficacy collapse.
\end{remark}

% =========================================================================
\section{Route Compression and Inoculum-Monotone Hitting Times}

\subsection{ODE system}

The standard target-cell-limited within-host model:
\begin{equation}
\frac{dT}{dt} = \lambda - dT - \beta TV, \quad
\frac{dI}{dt} = \beta TV - \delta I, \quad
\frac{dV}{dt} = pI - cV, \quad
\frac{dR}{dt} = \alpha I
\label{eq:ode}
\end{equation}
where $T$ = target cells, $I$ = infected cells, $V$ = free virions, $R$ = cells with stable proviral integration.

\begin{lemma}[Absorbing state]\label{lem:absorb}
The set $\mathcal{A} = \{R > 0\}$ is positively invariant. If $R(t_{0}) > 0$ then $R(t) > 0$ for all $t > t_{0}$.
\end{lemma}

\begin{proof}
$dR/dt = \alpha I \geq 0$, so $R$ is non-decreasing. Under ART, $I \to 0$ and $dR/dt \to 0$: the reservoir is static but does not shrink.
\end{proof}

\subsection{Route-dependent initial conditions}

$V^{(m)}(0) = 1$ virion/mL (mucosal, post-epithelial attenuation) versus $V^{(p)}(0) = 10^{3}$ virions/mL (parenteral, direct intravascular delivery). All ODE parameters in (\ref{eq:ode}) are identical across routes.

\begin{lemma}[Inoculum-monotone hitting times]\label{lem:monotone}
Solutions of (\ref{eq:ode}) with $V(0) = v_{0} < \tilde{V}(0) = \tilde{v}$ satisfy $\tilde{I}(t) \geq I(t)$ and $\tilde{R}(t) \geq R(t)$ for all $t \geq 0$. Therefore $\Tint(\tilde{v}) \leq \Tint(v_{0})$ a.s.
\end{lemma}

\begin{proof}
The subsystem $(I, V, R)$ is cooperative for fixed $T$:
\[
\frac{\partial(\beta TV)}{\partial V} \geq 0, \quad
\frac{\partial(pI)}{\partial I} \geq 0, \quad
\frac{\partial(\alpha I)}{\partial R} = 0.
\]
By comparison theorems for cooperative ODE systems\cite{smith1995}, solutions are order-preserving in initial conditions. Since $\tilde{V}(0) > V(0)$: $\tilde{V}(t) \geq V(t)$ and $\tilde{I}(t) \geq I(t)$ for all $t \geq 0$. Then
\[
\tilde{R}(t) = \alpha \int_{0}^{t} \tilde{I}(s)\,ds \;\geq\; \alpha \int_{0}^{t} I(s)\,ds = R(t),
\]
so $\Tint(\tilde{v}) \leq \Tint(v_{0})$.
\end{proof}

\subsection{Route compression corollary}

\begin{assumption}[Stochastic dominance]
$\Pint^{(p)}(t) \geq \Pint^{(m)}(t)$ for all $t \geq 0$. Follows from Lemma~\ref{lem:monotone} under route-specific inocula.
\end{assumption}

\begin{corollary}[Route compression]\label{cor:route}
Under Assumptions~A1--A4:
\begin{enumerate}[label=\textnormal{(\roman*)},leftmargin=2.5em]
    \item $\Epep^{(p)}(t) \leq \Epep^{(m)}(t)$;
    \item $\tcrit^{(p)}(\eta) \leq \tcrit^{(m)}(\eta)$.
\end{enumerate}
\end{corollary}

\begin{proof}
(i) The weight on $\emax$ equals $1 - \Pseed(t)$, which is smaller for parenteral; the weight on $\emin = 0$ equals $\Pint(t)$, which is larger for parenteral. So $\Epep^{(p)} \leq \Epep^{(m)}$.

(ii) Both functions are non-increasing; $\tcrit^{(p)}(\eta) = \inf\{t : \Epep^{(p)} < \eta\} \leq \inf\{t : \Epep^{(m)} < \eta\} = \tcrit^{(m)}(\eta)$.
\end{proof}

% =========================================================================
\section{Population-Level PEP Efficacy Bound}

\begin{definition}
$\Ebarpep = \mathbb{E}[\Epep(T_{\text{access}})] = \int_{0}^{\infty} \Epep(t)\,d\Faccess(t)$.
\end{definition}

\begin{proposition}[Population-level bound]\label{prop:bound}
Under Assumptions~A1--A4,
\begin{equation}
\Ebarpep \;\leq\; \Faccess(\tcrit)\,\emax + \big(1 - \Faccess(\tcrit)\big)\,\eta
\label{eq:bound_si}
\end{equation}
\end{proposition}

\begin{proof}
Partition the integral at $\tcrit(\eta)$:
\[
\Ebarpep = \int_{0}^{\tcrit} \Epep(t)\,d\Faccess + \int_{\tcrit}^{\infty} \Epep(t)\,d\Faccess
\;\leq\; \emax\,\Faccess(\tcrit) + \eta\,(1 - \Faccess(\tcrit)),
\]
using $\Epep(t) \leq \emax$ for $t \leq \tcrit$ and $\Epep(t) < \eta$ for $t > \tcrit$ (Theorem~\ref{thm:finite}(iv)).
\end{proof}

\subsection{Worked example (v2 verified values)}

Using the multiscale model output at the canonical operating point ($V_{0} = 10^{3}$ parenteral, CV $= 0.3$ on $\beta, c, \delta, \alpha$; commit \texttt{37e27ea}), $\tcrit^{(p)}(0.05) \approx 34.5$~h. Modeling access delay as log-normal (median 96~h, geometric SD 2.0):
\[
\Faccess(34.5~\text{h}) = \Phi\!\left(\frac{\ln 34.5 - \ln 96}{\ln 2}\right) = \Phi(-1.476) \approx 0.0699.
\]
With $\eta = 0.05$ and $\emax = 0.95$:
\[
\Ebarpep \;\leq\; 0.0699 \times 0.95 + 0.9301 \times 0.05 \;\approx\; 0.113.
\]
Envelope reading: expected population PEP efficacy below approximately 11.3\%.

\subsection{Sensitivity of the envelope to $\Faccess$ parameterization}

The median time-to-PEP-access for PWID is not directly anchored by primary timing data in the cited references (Taylor 2019\cite{taylor2019_si} is a 3-page commentary; Baugher 2025\cite{baugher2025_si} is NHBS-PWID surveillance focused on PrEP/syringe metrics, not PEP timing). The median is therefore a modeling choice informed by qualitative discussion of structural barriers in those references. Sensitivity:

\begin{table}[H]
\centering
\small
\caption{Envelope sensitivity to $\Faccess$ median (GSD $= 2.0$, $\tcrit^{(p)} = 34.5$~h, $\eta = 0.05$, $\emax = 0.95$).}
\begin{tabular}{@{}lcc@{}}
\toprule
\textbf{Median} & \textbf{$\Faccess(34.5\,\text{h})$} & \textbf{Envelope $\Ebarpep$} \\
\midrule
72~h (JID-consistent) & 0.121 & $\leq 15.9\%$ \\
96~h (canonical)      & 0.070 & $\leq 11.3\%$ \\
110~h                 & 0.054 & $\leq 9.8\%$  \\
120~h                 & 0.046 & $\leq 9.1\%$  \\
\bottomrule
\end{tabular}
\end{table}

The envelope bound is the most generous (least pessimistic) reading of population-level PEP efficacy under the framework; in independent cascade-based modeling of structural access (companion analysis; see main text Population-level efficacy section), point estimates of $P(R_{0}=0)$ for PWID under current US policy are 0.003--0.7\%, one to four orders of magnitude below the envelope. Both lines of evidence support the qualitative claim that PEP cannot serve as a population-level safety net for PWID under current structural conditions; the envelope provides the analytical bound, the cascade evidence provides the central estimates.

\subsection{Limitation: $\Faccess$ denominator}

$\Faccess$ derives from population-wide PEP-access surveillance, which includes exposures from both transmission-competent and virologically suppressed sources. Under U=U (Assumption~A4), exposures from suppressed sources are outside the framework's domain. Restricting $\Faccess$ to transmission-competent exposures only would reduce the relevant denominator and yield a tighter (smaller) $\Faccess$ and therefore a tighter envelope. The reported envelope is thus conservative with respect to U=U scoping: the true within-domain envelope is no larger than the reported value.

% =========================================================================
\section{Multiscale Within-Host Model: Implementation}

The multiscale within-host model combines stochastic founder-phase dynamics with deterministic ODE integration once productive infection is established. This section documents the implementation; all code is reproducible from \texttt{SRC/multiscale\_model/} at commit \texttt{37e27ea} of the project repository.

\subsection{Phase~1: Stochastic founder-phase tau-leaping}

At low founder population ($V_{0} \leq I_{\text{handoff}} = 10$), the dynamics are sensitive to extinction events; deterministic integration of (\ref{eq:ode}) systematically underestimates the extinction probability and overestimates the rate of reservoir establishment. Phase~1 uses tau-leaping integration of the discrete-state Markov chain with reactions:
\[
T \xrightarrow{\beta T V} I, \qquad I \xrightarrow{\delta} \emptyset, \qquad I \xrightarrow{p} I + V, \qquad V \xrightarrow{c} \emptyset, \qquad I \xrightarrow{\alpha} R.
\]
A fixed-delay buffer of $\tau_{\text{eclipse}} = 0.9$ days ($\approx 22$~h, Perelson 1996 Table~2)\cite{perelson1996_si} is enforced between cell infection and integration competence, ensuring that the lower bound on $\Tint$ matches the biological eclipse floor regardless of population dynamics.

A realization is classified as \emph{extinct} if both $I(t)$ and $V(t)$ reach zero before the population threshold $I_{\text{handoff}} = 10$ is exceeded. A realization is classified as \emph{integrated} if $R(t) \geq R^{*} = 10$ cells/mL is reached. For mucosal exposure ($V_{0} = 1$), extinction probability under CV $= 0.3$ is $\approx 68\%$; for parenteral exposure ($V_{0} = 10^{3}$, CV $= 0.3$), extinction is essentially zero ($<0.5\%$).

\subsection{Phase~2: Deterministic ODE handoff}

Once $I(t)$ exceeds $I_{\text{handoff}} = 10$, the dynamics transition to deterministic integration of (\ref{eq:ode}) using \texttt{scipy.integrate.solve\_ivp} with adaptive step size. The handoff continues until $R(t) \geq R^{*}$, at which point $\Tint$ is recorded.

\subsection{Heterogeneity}

Inter-individual variability is modeled as multiplicative log-normal noise on $\{\beta, c, \delta, \alpha\}$ with coefficient of variation $0.3$. The noise parameters are sampled once per realization; deterministic phase parameters use the sampled values throughout. The CV $= 0.3$ value is consistent with the standard within-host model heterogeneity assumption used in earlier published noise-anchored models.

\subsection{Sample size and seeding}

We use $N = 500$ realizations per (route, $V_{0}$, CV) cell. Random seeds are deterministic: \texttt{rng = np.random.default\_rng(base\_seed*1{,}000{,}000+k)} where $k$ is the realization index. This guarantees byte-identical reproducibility from a cloned repository at the cited SHA. Sensitivity to $N$ at the canonical operating point is reported in main text Methods; $\tcrit$ estimates are stable to $\pm 0.5$~h across $N \in \{200, 500, 1000, 5000\}$.

\subsection{Output: $\tcrit$ and $\Pint(t)$}

For each cell, the empirical CDF $\Pint(t)$ is constructed from the $N$-realization integration-time distribution. The route-specific efficacy curve $\Epep(t)$ is computed from $\Pint(t)$ via the master equation (\ref{eq:master_si}), and $\tcrit(\eta)$ is the smallest $t$ at which $\Epep(t) < \eta$. Reported $\tcrit$ values:

\begin{table}[H]
\centering
\small
\caption{Verified $\tcrit$ at $\eta = 0.05$ from \texttt{SRC/multiscale\_model/results\_v3/heterogeneity\_summary.csv}.}
\begin{tabular}{@{}llcc@{}}
\toprule
\textbf{Route} & \textbf{$V_{0}$} & \textbf{CV} & \textbf{$\tcrit(0.05)$ [h]} \\
\midrule
Parenteral  & $10^{3}$ & 0.0 & 32.5 \\
Parenteral  & $10^{3}$ & 0.3 & \textbf{34.5} \\
Parenteral  & $10^{3}$ & 0.5 & 36.0 \\
Mucosal     & 1        & 0.0 & 57.5 \\
Mucosal     & 1        & 0.3 & \textbf{60.5} \\
Mucosal     & 1        & 0.5 & 60.0 \\
\bottomrule
\end{tabular}
\end{table}

The compression ratio at the canonical operating point (CV $= 0.3$) is $60.5/34.5 \approx 1.75\times$.

% =========================================================================
\section{NHP Concordance --- Per-Timepoint Detail}

The multiscale model output was compared to NHP empirical protection rates at the timepoints reported in Tsai 1998\cite{tsai1995_si} and Otten 2000\cite{otten2000_si}. Concordance is defined as $|\Epep(\text{delay}) - \text{NHP protection}| < 0.20$ (within 20 percentage points).

\begin{table}[H]
\centering
\small
\caption{Per-timepoint NHP concordance at the canonical operating point (CV $= 0.3$). Reproducible from \texttt{SRC/multiscale\_model/results\_v3/nhp\_concordance.csv}.}
\begin{tabular}{@{}llccccc@{}}
\toprule
\textbf{Study} & \textbf{Route} & \textbf{Delay [h]} & \textbf{NHP protection} & \textbf{Model $\Epep$} & \textbf{$|$diff$|$ [pp]} & \textbf{Concordant} \\
\midrule
Otten 2000      & mucosal    & 12 & 1.00 & 0.95 & 5.0  & yes \\
Otten 2000      & mucosal    & 36 & 1.00 & 0.95 & 5.0  & yes \\
Otten 2000      & mucosal    & 72 & 0.50 & 0.016 & 48.4 & no  \\
Tsai 1998 IV    & parenteral & 24 & 1.00 & 0.50 & 50.0 & no  \\
Tsai 1998 IV    & parenteral & 48 & 0.50 & 0.00 & 50.0 & no  \\
\bottomrule
\end{tabular}
\end{table}

\textbf{Compression-ratio concordance.} The model's mucosal/parenteral $\tcrit$ ratio at $\eta = 0.05$ is $60.5/34.5 \approx 1.75\times$. The empirical NHP compression ratio between Tsai 1998 IV (parenteral half-protection at 48~h) and Otten 2000 (mucosal half-protection at 72~h) is $72/48 = 1.50\times$. The model's compression value is concordant with the empirical NHP range of 1.5--2$\times$.

\textbf{Per-timepoint absolute discordance.} The within-host kinetic model under-predicts protection at:
\begin{itemize}[nosep,leftmargin=2em]
    \item Tsai 1998 IV at 24~h (model 0.50 vs.\ NHP 1.00),
    \item Tsai 1998 IV at 48~h (model 0.00 vs.\ NHP 0.50),
    \item Otten 2000 at 72~h (model 0.016 vs.\ NHP 0.50).
\end{itemize}
We interpret this gap as evidence of NHP-specific protective contributions not captured by the within-host kinetic framework alone. Possible mechanisms include innate immunity contributions during the eclipse phase, mucosal containment by epithelial defense, and the small NHP cohort size ($n = 2$--4 per group) producing wider empirical confidence intervals on observed protection rates than the model's first-passage formulation accommodates. The compression-ratio result and the absolute-timepoint result are independent: the compression result is robust to the model's domain limitations; absolute-timepoint discordance constrains the interpretation of the model as a kinetic-only ceiling rather than a full mechanistic prediction of NHP protection.

% =========================================================================
\section{Parameter Tables}

\subsection{ODE system parameters (Perelson 1996 anchored)}

\begin{table}[H]
\centering
\small
\caption{ODE system parameters anchored on Perelson et al.\ 1996 (PMID 8599114).}
\begin{tabular}{@{}llllp{4.5cm}@{}}
\toprule
\textbf{Symbol} & \textbf{Description} & \textbf{Value} & \textbf{Units} & \textbf{Source} \\
\midrule
$\delta$ & Infected cell death rate & 0.7 & day$^{-1}$ ($t_{1/2} \approx 1.6$d) & Perelson 1996, Table~1 \\
$c$ & Viral clearance rate & 23 & day$^{-1}$ ($t_{1/2} = 0.24$d $\approx 6$h) & Perelson 1996, Table~1 \\
$\tau$ & Viral generation time & 2.6 & days ($\approx 62$h) & Perelson 1996, Table~2 \\
Eclipse phase & Min time to integration competence & 0.9 & days ($\approx 22$h) & Perelson 1996, Table~2 \\
$\beta$ & Infection rate constant & $2.4 \times 10^{-5}$ & mL$\cdot$virion$^{-1}\cdot$day$^{-1}$ & Standard model \\
$p$ & Virion production per infected cell & $10^{3}$ & virions/cell/day & Standard model \\
$\alpha$ & Stable integration rate & $10^{-3}$ & day$^{-1}$ & Calibrated \\
\bottomrule
\end{tabular}
\end{table}

Parameters $c$ and $\delta$ directly from Table~1; eclipse phase and generation time from Table~2. \textit{Note:} the reference $\log_{10}$ VL $= 4.5$ for PWID is from NHBS/NHAS surveillance, not from Perelson 1996 (whose patients had mean $\log_{10}$ VL $\approx 5.3$).

\subsection{Drug efficacy parameters}

\begin{table}[H]
\centering
\small
\caption{Stage-dependent drug efficacy parameters.}
\begin{tabular}{@{}llcp{6.5cm}@{}}
\toprule
\textbf{Parameter} & \textbf{Value} & \textbf{State $Z(t)$} & \textbf{Justification} \\
\midrule
$\emax$ (pre-seeding) & 0.95 & $Z(t) = 0$ & Near-complete suppression of initial RT; consistent with tenofovir/FTC NHP protection data \\
$\emid$ (seeded/pre-int) & 0.50 & $Z(t) = 1$ & Partial suppression of established replicative focus; sensitivity $\pm 40\%$ shifts $\tcrit < 6$h \\
$\emin$ (integrated) & 0 (fixed) & $Z(t) = 2$ & ART does not eradicate provirus; set to 0 by A2 and Lemma~\ref{lem:absorb} \\
\bottomrule
\end{tabular}
\end{table}

\subsection{Route-specific initial conditions and thresholds}

\begin{table}[H]
\centering
\small
\caption{Route-specific initial conditions and threshold parameters.}
\begin{tabular}{@{}lllp{6cm}@{}}
\toprule
\textbf{Parameter} & \textbf{Mucosal} & \textbf{Parenteral} & \textbf{Note} \\
\midrule
$V(0)$ virions/mL & 1 & $10^{3}$ & 3-log difference drives route compression; post-epithelial attenuation vs.\ direct IV delivery \\
$T(0)$ CD4 cells/mL & $10^{6}$ & $10^{6}$ & Identical; peripheral blood \\
$I(0)$ and $R(0)$ & 0 & 0 & No infected cells or reservoir at exposure \\
$I^{*}$ (seeding threshold) & 100 cells/mL & 100 cells/mL & Sensitivity: $[10, 1000]$; $\tcrit$ shifts $<8$h; compression ratio stable \\
$R^{*}$ (integration threshold) & 10 cells/mL & 10 cells/mL & Sensitivity: $[1, 100]$; parenteral/mucosal ratio stable 1.5--2$\times$ \\
$I_{\text{handoff}}$ (Phase 1$\to$2) & 10 cells/mL & 10 cells/mL & Stochastic-to-deterministic transition threshold \\
\bottomrule
\end{tabular}
\end{table}

\subsection{Multiscale simulation parameters}

\begin{table}[H]
\centering
\small
\caption{Multiscale simulation parameters.}
\begin{tabular}{@{}lll@{}}
\toprule
\textbf{Parameter} & \textbf{Value} & \textbf{Description} \\
\midrule
$N$ Monte Carlo replications & 500 / cell & Per (route, $V_{0}$, CV) cell; total $\sim$27{,}000 across full sweep \\
Heterogeneity CV & 0.3 & Multiplicative log-normal on $\{\beta, c, \delta, \alpha\}$ \\
Random seed & deterministic & \texttt{rng = np.random.default\_rng(42*1{,}000{,}000+k)} \\
Tau-leaping step & adaptive & Phase 1, founder-phase \\
ODE integrator & \texttt{solve\_ivp} (LSODA) & Phase 2, deterministic handoff \\
Eclipse delay buffer & 0.9 days ($\approx 22$h) & Fixed delay between cell infection and integration competence \\
CI type & 95\% credible interval & From stochastic $\Tint$ distribution \\
\bottomrule
\end{tabular}
\end{table}

% =========================================================================
\section{Stochastic VL Analysis --- Extended Methods}

\subsection{Source VL distributions}

Expected PEP efficacy was computed as a Monte Carlo integral over source viral load:
\begin{equation}
\mathbb{E}[\Epep(t)] = \int \Epep(t, \text{VL})\,P(\text{VL})\,d\text{VL}
\label{eq:mc_integral}
\end{equation}

\begin{table}[H]
\centering
\small
\caption{Source VL distributions parameterized from surveillance data. All log10-normal.}
\begin{tabular}{@{}llcp{6cm}@{}}
\toprule
\textbf{Distribution} & \textbf{Mean $\log_{10}$} & \textbf{SD} & \textbf{Empirical basis} \\
\midrule
PWID untreated/unsuppressed & 4.5 & 1.1 & NHBS/NHAS surveillance; chronic untreated HIV-1 in PWID networks \\
PWID mixed treatment & 3.2 & 1.5 & Mixture: suppressed ($\sim 1.5$) + unsuppressed ($\sim 4.5$) in partially treated PWID networks \\
General HIV+ community & 3.8 & 1.2 & Mixed treatment uptake; heterogeneous MSM/general community networks \\
Acute-infection-enriched & 5.0 & 0.9 & High-transmission networks; acute viremia peak; recent-infection clustering \\
\bottomrule
\end{tabular}
\end{table}

\textit{Note:} PWID untreated mean $\log_{10} = 4.5$ is from NHBS/NHAS surveillance, NOT from Perelson 1996.

\subsection{Three-regime temporal classification}

\begin{table}[H]
\centering
\small
\caption{Three-regime CI-width classification for the PWID untreated distribution.}
\begin{tabular}{@{}lp{4cm}lp{5cm}@{}}
\toprule
\textbf{Regime} & \textbf{Hours (PWID untreat.)} & \textbf{Peak CI width} & \textbf{Clinical interpretation} \\
\midrule
1 --- Early window & 0--22.0~h & $<$15~pp & VL uncertainty low-impact; initiate PEP immediately regardless of source VL \\
2 --- Critical window & 22.0--78.4~h (PWID/acute); 22.0--88.2~h (mixed/general) & 39.7~pp at 49~h & Maximal VL uncertainty; source VL testing clinically actionable; VL premium peaks at 21~h \\
3 --- Late window & 78.4--120~h (PWID/acute); 88.2--120~h (mixed/general) & $<$10~pp (collapsed) & Efficacy uniformly poor; VL irrelevant; $P$(futile) $> 90\%$ \\
\bottomrule
\end{tabular}
\end{table}

The Regime 1/2 boundary at 22.04~h is consistent across all four VL distributions (eclipse phase lower bound, Perelson 1996). The Regime 2/3 boundary is distribution-dependent: 78.37~h for PWID untreated and acute-infection-enriched distributions (higher VL accelerates CI collapse); 88.16~h for mixed-treatment and general-community distributions. This distribution-specificity is biologically coherent: higher inoculum populations exhaust the critical window earlier.

\subsection{VL knowledge premium and non-monotonic plateau}

VL knowledge premium $=$ absolute difference in expected $\Epep(t)$ between known vs.\ unknown source VL. Peak values: 3.76~pp at 21~h (acute-enriched); 2.94~pp (PWID untreated); 1.87~pp (general community); 1.27~pp (mixed Tx). All distributions peak at 21~h --- maximal VL-knowledge value at critical window onset.

% =========================================================================
\section{City-Stratified Analysis --- Extended Methods}

\subsection{Data source}

City-level HIV surveillance data were extracted from AIDSVu downloadable datasets for 34 high-burden US metropolitan areas (2023 reporting year; downloaded October 2025).

\subsection{Mixture-model VL derivation}

\begin{equation}
\mu_{\text{mix}} = f_{s}\,\mu_{\text{suppressed}} + (1 - f_{s})\,\mu_{\text{unsuppressed}}
\label{eq:mixture}
\end{equation}

Suppressed: $\log_{10}$ mean $= 1.5$, SD $= 0.4$. Unsuppressed: $\log_{10}$ mean $= 4.5$, SD $= 1.1$. IDU correction: $f_{s}$ reduced by 0.12~percentage points per percent IDU prevalence, capped at 12~percentage points total.

\subsection{Structural delay: composite indicator (not proxy)}

\begin{equation}
\Delta t_{\text{structural}} = \beta_{1}\,(\text{LateDx}\% - 10\%) + \beta_{2}\,\max(90\% - \text{Linkage}\%,\,0) + \Delta t_{\text{SDOH}}
\label{eq:delay_si}
\end{equation}

$\beta_{1} = 1.2$~h/\%; $\beta_{2} = 0.3$~h/\%. SDOH modifier 0--6~h from Gini + poverty. Barrier categories: low ($<$8~h), moderate (8--18~h), high (18--30~h), severe ($\geq$30~h).

\textbf{Methodological note.} City-level structural delay $\Delta t_{\text{structural}}$ is a composite indicator constructed as a calibrated linear function of two AIDSVu surveillance variables (late-diagnosis percentage; linkage-to-care percentage) plus an SDOH modifier. The composite is not a proxy for an unmeasured latent construct under proximal causal inference\cite{tchetgentchetgen2020_si}, and the analysis does not invoke the identification conditions required for proxy-based estimation\cite{pearl2009_si}. Coefficients $\beta_{1}, \beta_{2}$ are calibrated to reproduce the median city-level access friction; the composite is interpreted as a description of healthcare access friction at the city level, not as an estimand for a separate latent variable.

\subsection{Counterfactual analysis}

Parallel simulation assigning all 34 cities a uniform 24-hour access delay isolates structural barrier contribution from community VL. 33/34 cities show observed efficacy $>$ counterfactual (structural delays $<$24~h). Hartford CT is the sole exception (delay 24.4~h; high-barrier category, 18--30~h). Counterfactual range across cities: 76.0\%--82.1\%.

% =========================================================================
\section{Figure Inventory}

\textbf{Overview.} Three main figures (Figs.~1--3 of the main text), each up to four panels, generated by Python scripts in the project repository. All figures at 300~dpi. v2 Figures~1A and 2A include a vertical shaded band at $\log_{10}(V\!L) < \log_{10}(200)$ labeled "U=U region (per PARTNER2 / modern operational definition); model not applicable, transmission risk $\approx 0$". Figure~3 is unchanged: AIDSVu late-diagnosis and viral suppression data are aggregated across mode of HIV acquisition and cannot be reliably subset to U=U status.

\subsection{Figure 1 --- Route-Dependent PEP Efficacy Decay}

Source: \texttt{route\_models/make\_figures.py} (multiscale v3 outputs). Output: \texttt{Figure\_1\_Route\_Dependent\_PEP\_Efficacy\_Decay.png}.

\begin{table}[H]
\centering
\small
\begin{tabular}{@{}lp{4cm}p{8cm}@{}}
\toprule
\textbf{Panel} & \textbf{Title} & \textbf{Description} \\
\midrule
A & Mucosal vs.\ Parenteral Efficacy & $\Epep(t)$ vs.\ hours post-exposure for both routes at the canonical operating point ($V_{0} = 1$ mucosal, $V_{0} = 10^{3}$ parenteral, CV $= 0.3$). Vertical shaded band at $\log_{10}(V\!L) < \log_{10}(200)$ for U=U scoping. Triangles overlay NHP protection data. Primary panel. \\
B & State-Diagram and Inoculum Schematic & Three-state Markov diagram with route-dependent inoculum visualization. \\
\bottomrule
\end{tabular}
\end{table}

\subsection{Figure 2 --- Stochastic Efficacy Analysis Under Source VL Uncertainty}

Source: \texttt{stochastic\_layers/vl\_uncertainty.py} (multiscale-derived). Output: \texttt{Figure\_2\_Stochastic\_Efficacy\_VL\_Uncertainty.png}.

\begin{table}[H]
\centering
\small
\begin{tabular}{@{}lp{4cm}p{8cm}@{}}
\toprule
\textbf{Panel} & \textbf{Title} & \textbf{Description} \\
\midrule
A & Mean $\pm$ 95\% CI Efficacy Curves by Community VL Distribution & Mean $\Epep(t) \pm$ 95\% CI for all four VL distributions over 0--120~h. Vertical shaded band for U=U scoping. Regime shading overlaid. Primary panel. \\
B & Community VL Distributions & $\log_{10}$-normal density plots for all four VL distributions on common x-axis. \\
C & Probability of Sub-Therapeutic PEP by Delay and Community & $P(\text{efficacy} < 50\%)$ vs.\ hours post-exposure for all four distributions. Primary panel. \\
D & VL Knowledge Premium & VL knowledge premium (pp, absolute) vs.\ hours for all four distributions. Primary panel. \\
\bottomrule
\end{tabular}
\end{table}

\subsection{Figure 3 --- City-Stratified PEP Efficacy Across 34 US Metropolitan Areas}

Source: \texttt{stochastic\_layers/city\_stratified\_figures.py}. Outputs: \texttt{Figure\_3\_City\_Stratified\_PEP\_Efficacy.png} (4-panel, primary).

Figure~3 panels are unchanged from v1; the AIDSVu surveillance data underlying this figure is aggregated across mode of HIV acquisition and cannot be reliably subset to U=U status. The figure caption in the main text explicitly notes this scope limitation.

% =========================================================================
% Supplementary references
% =========================================================================
\renewcommand{\refname}{\large\textbf{Supplementary References}}

\begin{thebibliography}{9}
\bibitem{tanner2025_si}
M.R. Tanner \textit{et al.}, Antiretroviral postexposure prophylaxis --- CDC recommendations 2025. \textit{MMWR Recomm. Rep.} (2025).

\bibitem{siliciano2015_si}
J.D. Siliciano, R.F. Siliciano, Enhanced culture assay for detection of latently infected resting CD4+ T-cells. \textit{Methods Mol. Biol.} \textbf{1354}, 3--15 (2015).

\bibitem{wahl2023_si}
A. Wahl, L. Al-Harthi, HIV infection of non-classical cells in the brain. \textit{Retrovirology} \textbf{20}, 1 (2023).

\bibitem{perelson1997_si}
A.S. Perelson \textit{et al.}, Decay characteristics of HIV-1-infected compartments during combination therapy. \textit{Nature} \textbf{387}, 188--191 (1997).

\bibitem{perelson1996_si}
A.S. Perelson, A.U. Neumann, M. Markowitz, J.M. Leonard, D.D. Ho, HIV-1 dynamics \textit{in vivo}: virion clearance rate, infected cell life-span, and viral generation time. \textit{Science} \textbf{271}, 1582--1586 (1996).

\bibitem{mcmichael2010_si}
A.J. McMichael, S.L. Rowland-Jones, The immune response to HIV. \textit{Immunity} \textbf{33}, 431--441 (2010).

\bibitem{smith1995}
H.L. Smith, \textit{Monotone Dynamical Systems} (AMS Mathematical Surveys vol.~41, 1995).

\bibitem{tsai1995_si}
C.C. Tsai \textit{et al.}, Prevention of SIV infection in macaques by (R)-9-(2-phosphonylmethoxypropyl)adenine. \textit{Science} \textbf{270}, 1197--1199 (1995). [Timing data: \textit{J. Virol.} \textbf{72}, 4265--4273 (1998)]

\bibitem{otten2000_si}
R.A. Otten \textit{et al.}, Efficacy of postexposure prophylaxis after intravaginal exposure of pig-tailed macaques to a human-derived retrovirus (HIV type 2). \textit{J. Infect. Dis.} \textbf{182}, 1820--1825 (2000).

\bibitem{taylor2019_si}
J.L. Taylor, A.Y. Walley, A.R. Bazzi, Stuck in the window with you: HIV exposure prophylaxis in the highest risk people who inject drugs. \textit{Subst.\ Abus.} \textbf{40}, 441--443 (2019).

\bibitem{baugher2025_si}
A.R. Baugher \textit{et al.}, Are we ending the HIV epidemic among persons who inject drugs? Key findings from 19 US cities. \textit{AIDS} \textbf{39}(12), 1813--1819 (2025).

\bibitem{tchetgentchetgen2020_si}
E.J. Tchetgen Tchetgen, A. Ying, Y. Cui, X. Shi, W. Miao, An introduction to proximal causal inference. \textit{Statistical Science} \textbf{39}, 375--390 (2024).

\bibitem{pearl2009_si}
J. Pearl, \textit{Causality: Models, Reasoning, and Inference}, 2nd ed.\ (Cambridge University Press, 2009).
\end{thebibliography}

\end{document}
